% Lösungen Einheit 2 von 3 – Mit der Näherungsformel rechnen: die Grenzgleichung nach p lösen (quadratisch (zusammenbau v0.1)
\einheitenkopf{Einheit 2 von 3 · Mit der Näherungsformel rechnen: die Grenzgleichung nach p lösen (quadratisch}

\erg{11}{a) vorwärts: Grenzen gesucht \quad b) $[0{,}257; 0{,}347]$: die Grenzgleichung mit $h = 0{,}30$ und $c = 1{,}96$ nach p gelöst, beide Lösungen sind die Grenzen \quad c) 180 Treffer: $h = 0{,}319 + 1{,}96 \cdot \sqrt{0{,}319 \cdot 0{,}681 / 500} \approx 0{,}360$, erst h, dann mal 500 \quad d) $n = 246$: aus $0{,}05 > 1{,}96 \cdot \sqrt{0{,}2 \cdot 0{,}8 / n}$ folgt $n > 245{,}9$ – mit Anteilen rechnen, nicht mit Anzahlen}
\erg{12}{a) verträglich: $\mu = 56$, $\sigma \approx 4{,}10$, $\mu - 1{,}96\sigma \approx 47{,}97$, $\mu + 1{,}96\sigma \approx 64{,}03$; die Anzahl 50 liegt im Intervall}
\erg{13}{a) Richtig: $[0{,}251; 0{,}354]$ aus der Grenzgleichung, nach p gelöst. Fehler: Jan hat h statt p unter die Wurzel gesetzt \quad b) Nach dem Quadrieren entsteht eine quadratische Gleichung in p; die kleinere Lösung ist die untere, die größere die obere Grenze – zusammen begrenzen sie alle mit h verträglichen p}
